% =============================================================================
% NOTA DO ARQUIVISTA — Moldura Diegética do Arquivo 1.260
% =============================================================================
\cleardoublepage
\thispagestyle{empty}
\chapter*{Nota do Arquivista}
\addcontentsline{toc}{chapter}{Nota do Arquivista}
\markboth{Nota do Arquivista}{}

\begin{center}
\itshape\small
Este envelope contém a transcrição integral do diário do Paciente 1.260.\\
A leitura do material original causou alterações psicossomáticas nos examinadores.\\[0.5em]
\upshape\scriptsize --- Anotação em caneta-tinteiro no envelope pericial nº 1.260/1979
\end{center}

\vspace{1.2em}

\begingroup

O volume que você tem em mãos não foi escrito para ser publicado. É a reconstituição pericial do diário de \textbf{Joaquim Antônio Correia}, o Paciente 1.260 do Hospital Colônia de Barbacena.

Sobrevivente da Seca de 1915 e dos campos de concentração de retirantes no Ceará, Joaquim foi internado no Colônia em 1917 e lá permaneceu por 62 anos, até sua morte em 1979. Ao longo desse cárcere, produziu 847 páginas de anotações fragmentárias. O caderno original foi encontrado amarrado ao seu peito por tiras de pano quando seu corpo foi exumado da vala comum --- 62 dias após o sepultamento.

\vspace{0.4em}
\noindent Segundo o relato pericial do coveiro, o papel emanava um odor adocicado de umidade e terra, e a tinta parecia não secar. O laudo do IML registrou o achado como objeto de custódia e o anexou ao prontuário.

\vspace{0.6em}
Esta edição reúne 84 fragmentos organizados em cinco dimensões: as memórias da infância no sertão; o isolamento e as torturas no hospital psiquiátrico; a devassa dos laudos clínicos; a investigação forense conduzida pela Dra. Lúcia Mendes em 2026; e o registro da contaminação do próprio leitor.

\vspace{0.6em}
Não se trata de uma simples história de horror. A dor histórica que inspira estas páginas não é ficção. O prontuário de Joaquim, porém, é uma reconstituição ficcional; os números abaixo são estimativas atribuídas à bibliografia, não contagens deste arquivo.

\vspace{1.2em}
\noindent{\itshape
Recomenda-se não ler este livro no escuro.\\[0.35em]
Recomenda-se não ler este livro em solidão absoluta.\\[0.35em]
Se o cheiro de terra molhada surgir durante a leitura, feche o volume imediatamente.}

\vspace{1.2em}
\noindent A ferida histórica sobre a qual esta narrativa se assenta é real. O Hospital Colônia de Barbacena (1903--1994) internou à revelia cerca de 60 mil pessoas; 1.853 corpos foram exumados de suas valas comuns, segundo Daniela Arbex. Joaquim é um personagem, mas a dor que ele carrega --- fome, deslocamento, confinamento, tortura psiquiátrica --- pertenceu a milhares de brasileiros anônimos.

\vspace{0.6em}
\noindent O horror aqui evocado não deve ser consumido como entretenimento. A manipulação do leitor é uma escolha ética e formal --- não para enganar, mas para fazer o corpo sentir o que a mente sabe: que a vala não é uma metáfora, e que o sofrimento documentado existe para ser testemunhado, não estetizado.

\vspace{0.6em}
\noindent{\itshape Você será manipulado. Você será contaminado. Mas, ao final, saberá por quê.}

\vspace{1.4em}
\hfill{\scshape --- O Arquivista}\hspace{1em}

\endgroup
\clearpage
